\section{Related Work}

\paragraph{Attention complexity and FlashAttention.}
Self-attention scales quadratically in both time and memory with sequence length, which made long-context inference impractical for early Transformer models.
FlashAttention~\citep{dao2022flashattention} addressed the memory side by computing attention in tiles in on-chip SRAM, avoiding materializing the full attention matrix in HBM.
This brought memory use down to linear and cut HBM traffic substantially, making attention practical for sequences up to tens of thousands of tokens.
However, FlashAttention is exact (every attention score is still computed), so the quadratic time complexity remains.
As sequences grow attention computation becomes the dominant runtime cost and approximate methods become necessary.

\paragraph{Sparse attention for prefill.}
A natural way to reduce the time cost is to exploit the fact that attention matrices are empirically sparse: most tokens attend strongly to only a small fraction of others.
\citet{child2019generating} were among the first to demonstrate this, observing that trained Transformer attention matrices concentrate on local windows and a few distant positions.
They proposed fixed sparse patterns (local windows combined with strided or fixed summary positions) to approximate full attention with fewer FLOPs.
The limitation is that a fixed pattern cannot adapt to the content of a given input.

This motivated a shift to dynamic sparse attention, where the sparsity pattern is determined on the fly.
\citet{jiang2024minference} found that LLM attention heads exhibit consistent structural patterns across different inputs (A-shape, Vertical-Slash, and Block-Sparse) and assigned each head its dominant pattern offline, building sparse indices at inference time accordingly.
SpargeAttn~\citep{zhang2025spargeattn} is a parallel work that takes a more general approach designed to work across diverse model types beyond language models.
It selects a representative per block and uses representative-level attention scores to predict which block pairs are near-zero and can be skipped entirely, without assuming any fixed pattern.
FlexPrefill~\citep{lai2025flexprefill} and XAttention~\citep{xu2025xattention} are more recent works that push the state of the art further.
FlexPrefill adapts both the sparsity pattern and compute budget dynamically per head using Jensen-Shannon divergence to classify the pattern type at runtime.
XAttention uses the sum of antidiagonal values within each attention block as a lightweight proxy for block importance, achieving aggressive pruning at minimal overhead.

As a parallel line, \citet{qi2024deltallm}, a concurrent work from our group at LIACS, targets a different kind of sparsity.
While the methods above exploit the empirical sparsity of the attention map, they observe that adjacent key vectors are similar, so encoding each key as a delta from the previous one yields a sparse delta key matrix.
They then propose a method to compute attention directly using this delta key matrix, reducing the overall computation.
However, the delta is applied element-wise, giving unstructured sparsity scattered throughout the delta matrix, a pattern that GPU hardware cannot exploit efficiently, requiring custom hardware to accelerate.

\paragraph{KV cache compression for decoding.}
A separate line of work targets the memory bottleneck of autoregressive decoding.
Decoding is memory-bound: at every generation step the full KV cache must be loaded from HBM, so the KV cache dominates both memory capacity and memory bandwidth.
As context length grows, the cache can exceed GPU memory entirely, but even when it fits, a smaller cache directly translates to faster decoding by reducing the amount of data loaded per step.
\citet{zhang2023h2o} made the foundational observation that a small set of ``heavy hitter'' tokens receives the majority of attention mass, and proposed dynamically retaining only those tokens while evicting the rest.
\citet{li2024snapkv} refined this with a query-aware eviction policy: each attention head observes which tokens are attended to within an observation window at the end of the prompt, and retains only those positions for subsequent decoding.
\citet{zhang2024cam} pointed out that evicting tokens causes irreversible information loss, and proposed merging evicted tokens into their neighbors rather than discarding them.
\citet{wang2024kvmerger} arrived at merging from a different angle: they observed that adjacent token embeddings in the key matrix exhibit high cosine similarity within a sequence, and that this pattern is consistent across datasets and layers.
This motivates merging consecutive similar key tokens into a single representative, compressing the cache without throwing away information.

InfLLM~\citep{xiao2024infllm} spans both prefill and decoding.
It offloads the KV cache to CPU memory and selectively loads blocks to the GPU based on predicted attention relevance, using a block-representative scheme to decide which blocks a given query is likely to attend to.
Since CPU RAM is far more abundant than GPU VRAM, this effectively extends the usable context window without requiring the full cache on device at once.

\paragraph{Delta-based prefill acceleration.}
The three methods in this thesis are all training-free prefill acceleration methods built on a single observation: adjacent tokens have highly similar key and query vectors.
Each method applies this insight in a different setting, and each has a different closest relative among the methods above.

Delta-KV compression is a direct prefill optimization and is closest to KVMerger~\citep{wang2024kvmerger}.
Both compress the KV by merging adjacent keys that are nearly identical.
KVMerger does this to shrink the KV cache for decoding, while we compress $K$ at prefill to cut the cost of the $QK$ multiplication, turning the redundancy into structured sparsity along the sequence dimension.
To our knowledge this is the first work to use this redundancy as a prefill optimization with structured, GPU-friendly sparsity.

DeltaXAttention asks whether the delta principle can be layered on top of an existing block-sparse method rather than used on its own.
We take XAttention~\citep{xu2025xattention} and apply delta compression to the keys inside the tiles it selects, compressing within each tile on top of the tiles XAttention already skips.

Delta-based tile scoring (DPXA) is closest to XAttention~\citep{xu2025xattention} and SpargeAttn~\citep{zhang2025spargeattn}.
All three compute a cheap importance score for each tile to decide which tiles to compute exactly.
XAttention sums the antidiagonal values of each block, SpargeAttn mean-pools each block into a single representative, and DPXA compresses $Q$ and $K$ to form a proxy for the tile's attention mass.
